% Generated by scripts/build_textbook.py; edit content/ and rebuild.

\chapter{Open-Source \& Indigenous Digital Sovereignty}

\index{open source}

\index{CARE principles}
\index{community governance}
\index{data sovereignty}
\index{Indigenous data sovereignty}
\index{licence}
\index{Maori}
\index{OCAP}
Open-source software and Indigenous data sovereignty look like an odd pairing for a single chapter, until you notice that both make authority over digital material visible. Who decides how work or data may be used? Who benefits when it succeeds, and who carries the obligations of maintaining it? Open-source communities use copyright licences and governance structures to grant reuse rights and allocate decisions. Indigenous-led frameworks such as CARE, OCAP®, Māori data-sovereignty principles and Australian Indigenous data-sovereignty principles address rights and interests in data about Indigenous peoples, lands, resources and knowledge. \autocite{osi-definition,care,ocap,te-mana-raraunga,mnw-principles}


The comparison has strict limits. Mainstream open source generally begins with copyright held by creators and uses public licences to grant permissions. Indigenous data sovereignty begins from Indigenous peoples' rights and interests, including collective authority that conventional intellectual-property law may not recognise. The Indigenous-led frameworks named here also arise from different peoples and contexts; none substitutes for local authority. In the book's through-line, this is where professional practice becomes stewardship: the technical worker has to ask who has authority, who benefits, and what obligations survive after the system ships.


The chapter also refuses a second shortcut: no one Indigenous framework generalises everywhere. Māori, Aboriginal, Torres Strait Islander, First Nations and other Indigenous contexts differ in law, governance, history and protocol. Use the frameworks here as starting points for respectful professional questions, not as permission to proceed without local authority.


\index{community manager}
\index{community tech lead}
\index{FOSS}
\index{Open Source Program Office}
\index{te reo Maori}
\index{Te Hiku Media}
Across the sections that follow you will work through both toolkits: how to share data without surrendering guardianship, how projects govern themselves from solo stewardship to foundation backing, what a community tech-lead actually does all day, how FOSS licences allocate rights and obligations, and how Te Hiku Media built world-class speech technology while keeping te reo Māori under Māori control. Along the way, a running career thread (OSPO analysts, community managers, data kaitiaki, tech-leads) shows that stewarding the commons is not volunteer overflow work. It is a profession, and this part is your introduction to practising it well.

\section{Balancing Openness \& Cultural Safety}
\index{cultural safety}
\index{whakapapa}
Consider a composite scenario. A university language-technology lab asks to include a community organisation's archive in an open speech corpus. The collection includes recordings of elders who have since died, and the original permissions did not contemplate public release or AI training. The researchers see an opportunity to improve tools for an under-resourced language. The people connected to the recordings may see whakapapa, responsibilities, material that should not leave the community, and a history of research benefits flowing elsewhere.


The scenario has no culturally neutral ``open by default'' answer. The first professional task is to identify the people and institutions with authority to decide. The second is to separate possible purposes, users and derivatives rather than treating access as one switch. The third is to ask what benefit, capability and continuing control remain with the community.


\index{CARE principles}
\index{FAIR principles}
This is not an argument against sharing. The CARE Principles were developed because open-data and FAIR approaches do not by themselves address power differences, historical contexts, Indigenous rights and collective benefit. CARE complements technical data-management principles by centring people and purpose. \autocite{care}



\subsection{Use Frameworks In Their Proper Context}


Several Indigenous-led sources can improve the questions a team asks, but they are not interchangeable:


\begin{itemize}[leftmargin=*]
\item \textbf{CARE} stands for Collective benefit, Authority to control, Responsibility and Ethics. It is a global framework for Indigenous data governance and explicitly complements FAIR. \autocite{care}
\index{OCAP}
\item \textbf{OCAP®} stands for Ownership, Control, Access and Possession. It is specifically an expression of First Nations jurisdiction over information in Canada, not a generic set of principles for every Indigenous people. FNIGC also stresses that each First Nation may express the principles according to its own worldview and protocols. \autocite{ocap}
\index{data sovereignty}
\index{Indigenous data sovereignty}
\index{kaitiakitanga}
\index{Maori}
\item \textbf{Te Mana Raraunga's principles} articulate Māori data sovereignty through concepts including rangatiratanga, whakapapa, whanaungatanga, kotahitanga, manaakitanga and kaitiakitanga. \autocite{te-mana-raraunga}
\item \textbf{Maiam nayri Wingara's principles} provide an Australian Indigenous data-sovereignty source and should direct practitioners towards Indigenous peoples' right to exercise control over data ecosystems and their development. \autocite{mnw-principles}
\index{Traditional Knowledge}
\item \textbf{The United Nations Declaration on the Rights of Indigenous Peoples} provides an international rights context, including participation in decisions and control of cultural heritage, traditional knowledge and cultural expressions. It does not replace the authority or protocols of the community concerned. \autocite{undrip}
\end{itemize}

\index{tikanga}
Using a framework well begins with naming whose framework it is. A project involving Māori data should not cite OCAP® as if ``Indigenous'' were one jurisdiction. A Canadian project should not import tikanga terminology as decoration. A team in Australia should not treat national principles as permission to bypass the specific Aboriginal or Torres Strait Islander peoples whose data is involved.



\subsection{Consent Must Match The Proposed Use}
\index{consent}


A permission recorded for an oral-history project does not automatically cover public release, commercial use, voice cloning or model training. The difference is not hidden legal fine print; it changes who can hear the material, what can be inferred from it, and whether a derivative can be copied beyond recall.


Map the proposed uses separately:


\begin{itemize}[leftmargin=*]
\item listening or viewing by community members;
\item access by an approved researcher;
\item publication of selected excerpts;
\item creation of a searchable transcript;
\item training of a model;
\item release of model weights, embeddings or synthetic outputs;
\item commercial licensing or integration into another service.
\end{itemize}

For each use, record the relevant authority, purpose, users, duration, benefit, review point and withdrawal or suspension process. If the people with authority distinguish public, community-restricted and closed material, the system must preserve those distinctions. The technical team does not invent the categories.


Provenance must travel with the material. A future operator needs more than a file name and copyright field: they need to know which community or people are connected to it, the basis of authority, approved purposes, restrictions on derivatives, and who can answer questions. Local Contexts provides Traditional Knowledge and Biocultural Labels and related notices intended to make community protocols and provenance visible in collections and data systems. Teams should follow Local Contexts' process rather than assigning labels themselves. \autocite{local-contexts-labels}



\subsection{Turn Authority Into System Behaviour}


Governance has to survive contact with software. The appropriate controls depend on decisions made by the relevant authority, but a technical design may need:


\begin{itemize}[leftmargin=*]
\item access groups that reflect approved roles and purposes;
\item separate permissions for viewing, downloading and creating derivatives;
\index{renewal}
\item time-limited access with review before renewal;
\item logs that community decision-makers can inspect and understand;
\item encryption and key-management arrangements that do not leave control solely with an external partner;
\item alerts and suspension paths for use outside an approved plan;
\index{backups}
\item metadata that remains attached through export and backup;
\item documented deletion, return or preservation procedures when a partnership ends.
\end{itemize}

These are implementation options, not a universal cultural-safety checklist. A technically strict control can still be wrong if the wrong institution designed it or if it prevents the community from exercising its own rights. Conversely, a memorandum of understanding without enforceable access paths may express good intentions while leaving the platform's default permissions in charge.



\subsection{Agreements Need Benefit, Review And Exit}


A data-sharing agreement should be specific enough to operate. It can identify approved purposes, decision-makers, prohibited uses, reporting obligations, benefit-sharing, cultural briefings, publication review, incident handling and dispute resolution. It should also explain what happens to copies and derivatives when approval expires or a project ends.


Build review into the relationship. New model capabilities, a proposed commercial partner or a change in community priorities may make an earlier decision obsolete. The agreement should say who reconvenes, what evidence they receive and whether use pauses during review. A sunset date can force a decision that would otherwise drift into permanent access.


\index{licence}
Labels, licences and contracts support governance; none creates authority by itself. The people making decisions need time, resources, technical explanations they can challenge, and the ability to say no without losing unrelated services or funding.



\subsection{A Defensible Outcome For The Scenario}


The lab in the opening scenario cannot design the outcome on the community's behalf. One possible process is to return with a plain-language proposal that separates research access, public release and model training; resource the community's review; and accept that some material may remain closed. The decision-making authority may approve a limited pilot, require changes, postpone it or decline it.


If a pilot proceeds, its conditions might include time-limited access, no redistribution, agreed outputs, benefit returned in a useful form, and review before derivatives are retained or released. Those terms are examples for negotiation, not recommendations that override local protocols.


The professional habit is to replace ``can we make this open?'' with a set of accountable questions: who has authority, which use is proposed, who benefits, what obligations travel with the material, and how can the decision be changed? The answer may include openness, but openness is an outcome of governance rather than a default set by the person holding the file.

\section{Governance Journeys}
\index{maintainer}
\index{open source}
Much open-source infrastructure begins with one person sharing a useful tool, then acquires users and organisational dependants faster than it acquires maintainers. Governance is how a community decides who may act, who has a vote, how conflict is handled and how money, if there is any, gets spent.


The useful mental model is progressive scaffolding. You don't pour concrete foundations for a garden shed, but you also don't balance a skyscraper on a folding chair. Projects rarely jump from hobby repo to incorporated nonprofit overnight; healthy governance grows in layers. Early contributors document norms alongside code. Mid-stage teams add facilitation rituals. Mature foundations balance budgets, legal protection and public accountability. The test for each new layer is the same: it should lower the load on individuals and widen community agency, not calcify power around whoever got there first. And because contributors span cultures and time zones, governance is people-work — designing decision paths that are legible to newcomers, respectful of marginalised voices, and able to change when circumstances do.


\index{Te Hiku Media}
Where a project's community is grounded in its own cultural governance, the structures may look different. Te Hiku Media, examined later in this part, is one such case. Treat what follows as a mainstream open-source toolkit to compare with other systems, not as a universal model.



\subsection{Solo stewardship}


In the single-maintainer phase, governance is basically the maintainer's judgement plus whatever norms they capture in the README. Decisions are fast, vision is coherent, and the whole thing is one bad month from collapse. The jargon for this is the \emph{bus factor}: how many people can get hit by a bus before the project dies. Aim higher than one.


Solo governance still deserves artefacts. A CONTRIBUTING.md that sets expectations, a regular issue-triage rhythm, and an explicit roadmap (including what will \emph{not} be built) let the community help instead of guess. The risks are predictable: burnout, security fixes that ship unreviewed because there is nobody to review them, and knowledge that lives in one head. The mitigations are equally predictable and routinely skipped. Recruit trusted lieutenants and give them triage or release permissions, even on a trial basis; the people to trust are the ones who already follow the contribution guidelines reliably, communicate edge cases, and respect boundaries. Delegate release tokens. Document the deployment passwords. Schedule real vacations before the 3am ``the server is down and only I know the password'' stress arrives, because it will.



\subsection{Core team collectives}


Once three to ten maintainers coordinate, clarity beats charisma. This is the stage for onboarding playbooks that explain code-review expectations, decision timelines and how to escalate conflict — and for rotating roles like review captain, release manager and community moderator so knowledge spreads and nobody becomes the permanent bottleneck.


Mastodon is the instructive example. Before its restructure, founder Eugen Rochko was personally handling code reviews, harassment reports and server bills — no sustainable separation of duties, everything routed through late-night direct messages. The shift created squads for the iOS and Android clients, server administration and community moderation, each with two or three maintainers, plus shared decision logs in public forums.


Two practices carry most of the weight at this stage. The first is \emph{lazy consensus}: a proposal passes unless someone objects within a set timeframe, with a fallback vote if objections can't be resolved. It keeps decisions moving without demanding a meeting for every change. The second is transparency: publish decision logs somewhere public (GitHub Discussions, an open Notion space) so contributors can follow the reasoning rather than reverse-engineer it. Pair both with mentorship cohorts aimed at underrepresented regions, so contributors in Latin America or South Asia know how to surface blockers despite the time zone gap.



\subsection{Foundations and fiscal sponsors}


When a project becomes critical infrastructure, a foundation can add legal, fiscal and reputational scaffolding. Foundations and fiscal sponsors may hold trademarks, accept funds, contract staff, arrange insurance and support cross-project work. Those services and governance models change, so a project should compare current charters, fees and decision rights before choosing a home.


\index{licence}
Expect the paperwork to harden: contributor licence agreements, governance charters and codes of conduct become formal documents rather than wiki pages, and a board plus a technical steering committee keeps roadmap authority with engineers while adding accountability and succession planning. That bureaucracy buys real things — vendor-neutral roadmaps, multi-year funding commitments, and a buffer when commercial or geopolitical interests collide with community priorities, which they eventually do. If all you actually need is a bank account, a fiscal sponsor such as Open Collective Foundation can handle the money without the full institutional apparatus.



\subsection{When to formalise — and how it goes wrong}


\index{community governance}
\index{RFC}
\index{RFC process}
The trigger for adding structure is pain, and the pain has leading indicators: security patches ageing past thirty days, Fortune 100 adopters asking for SLA-like assurances, maintainers skipping parental leave because nobody else can cut a release. A short community health survey surfaces problems before they become departures. Two questions do a lot of work: ``How long do security fixes usually take?'' and ``Do you feel comfortable challenging technical decisions?'' Ask them regularly and track the trend. React took three years to move from a Facebook-internal tool to a more open governance model as enterprises demanded neutrality; the hand-off included an RFC process and a cross-company steering group with final authority. Whatever you choose, publish the timelines, the decision criteria and the retrospectives, so stakeholders understand why this structure matches the current risk.


Formalising badly has its own catalogue of failure modes:


\begin{itemize}[leftmargin=*]
\index{BDFL}
\item \textbf{BDFL burnout.} A founder clings to every decision until they crack. Rotate authority and document delegation triggers before the crisis.
\item \textbf{Committee paralysis.} Every choice needs consensus from a dozen people, so nothing ships. Set quorum rules, empower working groups, time-box debates.
\item \textbf{Corporate capture.} One vendor supplies enough funding or staff to bend the roadmap. Diversify revenue, publish conflict-of-interest disclosures, and design board representation and voting rules before dependence becomes control.
\item \textbf{Fork wars.} Communication breaks down and disagreement turns personal. Invest in mediation channels, transparent decision logs and cultural competency training so disputes stay technical.
\end{itemize}

Run incident reviews on governance failures the way you would on outages: the point is to learn, not to repeat the cycle with new personnel.



\subsection{The playbook, and the people who run it}


Good governance is documented well enough that newcomers can self-serve. Concretely: a GOVERNANCE.md describing decision rights, a CODEOWNERS file routing reviews, and an RFC template that shows a proposal's stages. Add onboarding kits with buddy assignments, office hours in multiple languages, and primers on inclusive communication. Track community health like you'd track uptime (median PR review time, the ratio of first-time contributors whose work gets merged, moderation responses within twenty-four hours) and share the dashboards publicly. Match decision frameworks to the decision: consensus-seeking for technical design, majority vote for budget approvals, veto powers reserved for safety issues. Pair transparent funding reports with conflict-resolution channels facilitated by trained moderators or ombudspeople.


\index{community manager}
\index{data sovereignty}
\index{Indigenous data sovereignty}
\index{Open Source Program Office}
All of this is paid work as well as volunteer work, which makes it a career path. Sustainable projects blend maintainers, release engineers, program managers, community stewards, translators, legal counsel and accessibility reviewers, with roughly one community manager per two hundred active contributors — reach that headcount and your ``side project'' problems have officially become good problems. Map the geography too: leadership shouldn't be North America-only, budgets should include stipends for maintainers in underrepresented regions, and meetings should rotate time zones. People arrive via contributor streaks, Google Summer of Code, corporate OSPO rotations and fellowships centring historically excluded communities. The ones who thrive communicate transparently, mediate cross-cultural tension, and respect Indigenous data sovereignty where projects touch it. From maintainer, the arc runs to technical steering chair, then foundation executive or policy advisor.



\begin{quote}\small
Rule of thumb: if a governance question keeps getting answered in private messages, it isn't governed yet. Write it down where the next person can find it.
\end{quote}


The takeaway: governance is not paperwork bolted onto code, it is the mechanism by which a project scales participation without losing its values. Intentional structural choices (paired with documentation, mentorship and cultural humility) are what separate the projects that outlive their founders from the ones that burn them out.

\section{Community Tech-Leads}
\index{community tech lead}
\index{maintainer}
\index{open source}
\index{RFC}
\index{RFC process}
Sarah's startup (last seen hiring in Part 6) builds on an open-source web framework, and one of her developers dutifully upstreamed a fix for a bug the team had hit in production. The pull request sat for five weeks. The RFC thread that might have made the fix unnecessary had two hundred comments and no decision. Nobody was being lazy; the project had simply grown past the point where a loose collection of maintainers could watch the whole ecosystem. Then the project's fiscal sponsor funded a community tech-lead role, and within a quarter the merge queue moved, RFCs had named facilitators, and Sarah's developer got actionable review feedback in three days instead of thirty-five. From the outside it looked like magic. From the inside it was a job — and this section is about that job.



\subsection{Why the role exists}


\index{localisation}
Community tech-leads are connective tissue: they sit between governance, engineering and contributor care, in the gap between maintainers, contributors and users that opens up once a project matures. Their core trick is translation — turning community values into technical direction and roadmap trade-offs, so the governance charter and the codebase don't drift apart. They also coordinate the work that no single maintainer can own: security response, accessibility, localisation.


There is no universal contributor-to-lead ratio. The evidence for funding the role should come from the project's own queue: review latency, unanswered proposals, moderation escalations, maintainer load and contributor departures. Without enough accountable leadership, RFCs stall, governance documents drift away from practice and trust erodes one unanswered contribution at a time.



\subsection{What a community tech-lead actually owns}


\index{SBOM}
The responsibilities split into facilitation and stewardship, and the facilitation half is not the soft half. Tech-leads run backlog triage, shepherd RFC discussions and coordinate release planning — always with transparent decision logs, so contributors can follow the reasoning without having been in the room. They model inclusive communication and uphold the code of conduct in technical forums, which is where codes of conduct most often quietly die. They pair emerging maintainers with mentors and steward succession plans for critical subsystems, so no module has a bus factor of one. On the stewardship side they own dependency hygiene, keep the SBOM current (the licensing section later in this part explains why auditors care), and coordinate incidents across distributed teams. And they report progress to fiscal sponsors, foundations and partner organisations — the role is accountable to people, not just to code.


\index{escalation path}
\index{licence}
Then there is the bridging work, which is the part most engineering roles never see. A tech-lead hosts contributor office hours across time zones and languages. They convert user research, localisation feedback and Indigenous data protocols (the kind covered earlier in this part) into actionable engineering issues, so cultural requirements land in the sprint rather than in a slide deck. They maintain the unglamorous tooling that keeps a community functioning: translation pipelines, documentation builds, contributor analytics. They negotiate with corporate stakeholders so funding aligns with the community roadmap instead of quietly redirecting it. And they act as the escalation path for moderation teams when a dispute turns out to have technical roots — a licence disagreement dressed up as a flame war, say.



\subsection{The shape of a day}


Because contributors are everywhere, the day is structured around time zones rather than a single office. Mornings go to async stand-up posts, merge-queue review and security patch triage — the things that block other people if they wait. Midday brings the roadmap sync with the governance committee, unblocking volunteer maintainers, and updating the mentoring roster. Evenings often hold a community livestream or a timezone-friendly pairing session, because ``the community'' includes the people for whom your morning is the middle of the night.


The allocation varies with the project's stage, but technical design is only one part of the job. Community facilitation, mentoring, moderation and reporting can occupy most of a week. A useful service-level measure is the proportion of contributions that receive actionable feedback within the published response window. They need not be merged; they do need an accountable answer.



\subsection{Pathways in, and the traits that stick}


\index{iwi}
There is no mysterious anointing. The most common route is the long-term contributor promoted after stewarding a module through multiple releases — the promotion formalises trust that already exists. Fellowship programs like Outreachy, Google Season of Docs and GitHub's OSS maintainer scholarships feed the pipeline. Engineers rotate out of corporate OSPOs looking for community-facing leadership. And community organisers come in the other direction, adding technical depth through bootcamps or iwi digital innovation hubs — a pathway that matters, because it says plainly that Indigenous technologists belong in these roles. The typical baseline is four to six years of combined engineering and facilitation experience with demonstrated community impact — a bar about skill depth, not about which degree you hold.


\index{blameless post-mortems}
\index{cultural safety}
The traits that keep people effective in the seat are consistent. Bilingual or multilingual communication, because global communities don't all argue in English. High empathy and conflict navigation grounded in cultural safety principles. Systems thinking — the ability to see how infrastructure, funding and governance intersect, so you fix causes rather than symptoms. Calm under incident pressure, including the ability to run blameless retrospectives that include community voices, not just staff. And a documentation-first mindset, so every decision is legible to the newcomer who arrives six months later.



\begin{quote}\small
A good tech-lead's fingerprints are hardest to see in the moment: the flame war that didn't happen, the maintainer who didn't burn out, the RFC that closed in three weeks instead of a year.
\end{quote}



\subsection{The team around the role, and where it leads}


\index{community manager}
Nobody does this solo. Tech-leads partner closely with product and community managers, legal counsel and accessibility reviewers, and they oversee the volunteer squads that do the recurring work: docs, localisation, security response, release engineering. A typical staffing shape is one community tech-lead supported by two or three senior maintainers and four to six rotating fellows. Part of the job is advocacy inside that structure — arguing for stipend budgets and equitable recognition programs, and making sure foundation boards hear frontline contributor data before they vote on policy, rather than after.


\index{data sovereignty}
\index{Indigenous data sovereignty}
\index{Open Source Program Office}
\index{runway}
Nor is it a cul-de-sac. The lattice runs from maintainer or module steward, through community tech-lead accountable for cross-project initiatives and contributor health metrics, up to ecosystem lead or OSPO director guiding multi-project strategy and funding partnerships. Executive pathways include foundation CTO, cooperative digital steward, and policy advisor on open-source sustainability, with lateral moves into developer relations, governance, or Indigenous data sovereignty leadership. For a role that many organisations still forget to fund, it has an unusually long runway.


The takeaway is the one Sarah's team learned from the outside: intentional investment in community tech-leads is what sustains contributor trust, equitable governance and long-term project resilience. Funding the role isn't adding a title — it protects the project's values, speeds its reviews, and keeps its governance tethered to the community it claims to serve. If you find yourself contributing to a project where reviews take five weeks and nobody owns the gap, you now know both the diagnosis and the job description for the cure.

\section{FOSS Licensing Choices}
\index{FOSS}
\index{copyleft}
\index{licence}
Six weeks into the hiring spree we watched in Part 6, one of Sarah's new developers needs to parse dates in three formats. He finds a tidy JavaScript helper on someone's blog, pastes it into the product, and moves on. There's no licence file and no attribution — which does not mean the code is free to use. With no permission, Sarah's company may be shipping code it has no right to reproduce or distribute. A copyleft notice would create a different analysis: the obligations depend on the licence, what has been modified or combined, and how the product is conveyed or made available. Either problem tends to surface at a bad moment, such as an automated scan during investor due diligence.


\index{open source}
Free and open-source software (FOSS) licences make those permissions explicit. The Open Source Definition describes criteria including free redistribution, source-code availability and permission for derived works; individual licence texts then state the operative rights, conditions and disclaimers. Read the actual text and obtain legal advice for material risk. \autocite{osi-definition}



\subsection{Permissive licences: MIT, BSD and Apache 2.0}
\index{Apache License}
\index{BSD License}
\index{MIT License}
\index{permissive licence}


Permissive licences generally grant broad rights to use, modify and redistribute code, including in closed-source products, while requiring preservation of notices or licence text. They commonly disclaim warranty and liability; the exact conditions differ, so do not substitute a category label for the licence.


MIT and BSD are the shortest and loosest. Apache 2.0 adds two things enterprise risk teams care about. First, an explicit patent grant: every contributor licenses any patents their contribution would otherwise infringe. Second, a retaliation clause: if you sue anyone claiming the project infringes your patents, your own patent licence to the project terminates. That combination is why large companies often prefer Apache 2.0 over MIT for anything patent-adjacent.


Because permissive licences allow proprietary derivatives, they're the default choice for start-ups building commercial services, vendor integration teams, and agencies embedding open libraries in client work. React.js is the canonical example: Facebook shipped it under MIT so that product teams anywhere could build proprietary apps on top without asking permission — and adoption exploded accordingly.


\index{maintainer}
The trade-off is reciprocity. A permissive licence gives you no legal lever to force improvements back upstream. Sustaining the project then depends on community goodwill, healthy governance, or a parallel commercial offering. If you're a first-time maintainer staring at the fine print, choosealicense.com and GitHub's built-in licence picker will walk you through the options before you click ``public''.



\subsection{Copyleft: the GPL family}
\index{GPL}


\index{AGPL}
\index{LGPL}
Copyleft licences (including GPLv2, GPLv3, LGPL and AGPL) attach source-code and licensing conditions to specified forms of copying, modification, conveying or network use. Whether those conditions extend to a combined work is a legal and technical question about the licence and the way components interact, not a rule that any nearby proprietary code automatically becomes open source. \autocite{gnu-licenses}


The mechanics matter, so be precise about them:


\begin{itemize}[leftmargin=*]
\item GPL source obligations are generally tied to conveying covered or modified work rather than purely private use. The form of delivery and the work actually conveyed matter.
\item \textbf{LGPL} is the softer variant for libraries: you may link it into a proprietary application without relicensing your own code, provided users can swap in and relink updated versions of the library.
\item \textbf{AGPL} adds a network-interaction provision: when users interact remotely with a modified covered programme, the licence requires an opportunity to receive the corresponding source. It does not simply redefine every network interaction as distribution.
\end{itemize}

The Linux kernel is GPLv2, which is why Android handset makers must publish their kernel modifications every time they release a firmware image — some learned this the hard way. Same story if you ship a router with modified GPL firmware: matching source must be made available. Copyleft rewards collaboration-heavy ecosystems — public-sector platforms, civic tech, and organisations like Signal, which keeps its core GPL precisely so its privacy claims stay independently inspectable.



\subsection{Mixing licences without getting burned}


Real products combine dozens of dependencies, and not all licences coexist happily. The working rules of thumb:


\begin{itemize}[leftmargin=*]
\item \textbf{MIT/BSD} combines with almost anything; just preserve attribution.
\item \textbf{Apache 2.0} combines with MIT, other Apache code and GPLv3 — but not GPLv2, whose terms clash with Apache's patent provisions. (GPLv3 was written partly to fix this.)
\item \textbf{GPLv3} code can absorb permissive code, but the combined binary you distribute must be GPL.
\item \textbf{AGPL} requires special attention for modified covered programmes offered over a network; read the licence rather than relying on a compatibility slogan.
\end{itemize}

Use a maintained compatibility source and legal review for consequential combinations so nobody discovers a licence conflict at the eleventh hour.


Some projects deliberately run two licence streams. MySQL and Qt offer their code under the GPL for the community and sell commercial licences to companies that want proprietary terms — a model that funds development while protecting openness goals. It only works if the steward holds the rights to relicense (more on that below) and is transparent about which features sit under which stream. MongoDB's shift to the Server Side Public License (SSPL) shows the same lever pulled defensively: a licence change designed to protect a business model from cloud providers reselling the product.


At the other extreme, CC0 lets a rights-holder waive copyright and related rights to the extent legally possible and supplies a fallback public licence where the waiver is ineffective. That structure accounts for differences between jurisdictions; it is not a declaration that every kind of right can always be abandoned. \autocite{cc0}


Copyright terms, moral rights, database rights, Crown copyright and patent scope differ across jurisdictions. Established licence texts reduce avoidable ambiguity, but OSI approval is not a substitute for jurisdiction-specific advice. Document the chosen licence, notices, provenance and any legal review.



\subsection{Contributor agreements and choosing for your context}


\index{DCO}
Once outsiders contribute, rights may be held by multiple people or employers. Contributor Licence Agreements can grant a project specified rights in contributions; their scope varies and should be read. Distinguish individual from corporate agreements, and check that the signatory has authority. Some projects instead use the Developer Certificate of Origin, under which contributors certify statements about the origin of a contribution by adding a \texttt{Signed-off-by} line. A DCO is a certification mechanism, not a copyright assignment. \autocite{dco}


\index{Open Source Program Office}
Choosing a licence is a cross-functional decision, not a checkbox. Start from goals: are you optimising for adoption, reciprocity, revenue, or community trust? The quick flow — need maximum adoption, go permissive; need guaranteed sharing, pick copyleft; need revenue plus openness, explore dual licensing. A worked example: a government agency building a records platform wants vendor fixes funded by taxpayers to stay public, so its open-source program office (OSPO) recommends GPLv3, legal sets up DCO-plus-CLA contribution intake, and the comms team briefs suppliers on what they're signing up for. Documenting that rationale now prevents a heated argument three years later when a new contractor joins. Involve legal counsel, community tech-leads, security engineers and product managers early — a licence choice is like a roommate agreement, and if you skip the awkward conversation about chores and guests, you'll have it later, angrier, with the mess already on the floor.



\subsection{Staying compliant — and who gets paid to care}


\index{containers}
\index{SaaS}
\index{SBOM}
Compliance is a process, not an event. Maintain a software bill of materials (SBOM) so you can prove which licences are in your build, and wire licence scanners into CI so incompatible combinations fail fast rather than surface in an audit. Record every redistribution moment (shipped binaries, published container images, a SaaS endpoint incorporating AGPL code) and keep source archives, NOTICE files and third-party attributions ready to go. Tools like FOSSA, OSS Review Toolkit and GitHub's dependency review do the tedious parts once configured. Companies have paid real settlements for ignoring GPL notices, so treat ``we'll fix the licensing post-launch'' as the red flag it is.



\begin{quote}\small
Licence compliance is like flossing: tedious, easy to skip while everything seems fine, and the neglect only announces itself as an expensive, painful audit.
\end{quote}


This work is a career path as well as a responsibility distributed across engineering, legal, security and procurement. OSPO analysts may arrive from software engineering, developer relations, compliance or technology law. The work rewards careful evidence, diplomacy and the ability to translate legal nuance for technologists without flattening it. Senior roles coordinate licence strategy, contribution policy and ecosystem relationships across business units.


\index{data sovereignty}
\index{Indigenous data sovereignty}
The takeaway: licence selection is a strategic lever, not paperwork. It signals how you invite collaboration, how you protect contributors, and (when projects touch cultural knowledge, as later sections of this part explore) how you honour data sovereignty commitments. Teams that understand reciprocity obligations, international quirks, CLAs and dual licensing can design contribution models that sustain trust with communities, customers and partners for the long haul.

\section{Te Hiku Media Case Study}
\index{Te Hiku Media}
\index{data sovereignty}
\index{Indigenous data sovereignty}
The previous sections argued that governance decides who benefits from digital resources. Te Hiku Media provides a public example of Māori-led media and language technology in which data sovereignty is part of the technical work rather than a policy added afterwards.


\index{kaitiakitanga}
\index{licence}
\index{Maori}
\index{te reo Maori}
This case needs a stricter reading discipline than a fictional business scenario. Te Hiku Media speaks for its own work; this book should not fill gaps in that account with plausible-sounding details about governance boards, cloud regions, contracts or cultural practice. The public evidence establishes that Te Hiku operates media and technology initiatives for te reo Māori, presents Papa Reo as a language platform for a multilingual Aotearoa, publishes work on te reo Māori speech recognition and synthetic voice, and provides resources on data sovereignty and its Kaitiakitanga License. The organisation's site also states that its hosted content is protected under that licence. \autocite{te-hiku-tech,te-hiku-kaitiakitanga}


Those facts are enough to make the professional point. A community can be the designer and operator of language technology, not merely a source of recordings for somebody else's system. It can also state terms for how language material is used rather than accepting ``open'' as the only legitimate technical default.



\subsection{Read The Case Through Māori Data-Sovereignty Principles}


\index{whakapapa}
Te Mana Raraunga's Māori Data Sovereignty Principles provide a stronger frame than a generic privacy checklist. They address rangatiratanga, whakapapa, whanaungatanga, kotahitanga, manaakitanga and kaitiakitanga in relation to Māori data. The principles concern authority, relationships, collective benefit, care and stewardship; they cannot be reduced to where a server is located. \autocite{te-mana-raraunga}


That distinction changes project questions. A conventional technology assessment may ask:


\begin{itemize}[leftmargin=*]
\item Is the dataset lawfully collected?
\item Is storage encrypted?
\item Does the model meet an accuracy threshold?
\item Can researchers reproduce the result?
\end{itemize}

A data-sovereignty assessment must also ask:


\begin{itemize}[leftmargin=*]
\item Who has authority to define the purpose of the work?
\item What relationships and whakapapa attach to the material?
\item Who controls later uses, including model training and commercial reuse?
\item How does the work return benefit and capability to Māori?
\item Who carries kaitiaki responsibilities through the system's life?
\end{itemize}

These questions do not supply a universal answer. They identify decisions that cannot be made by an external engineer merely because the engineer controls the infrastructure.



\subsection{A Licence Is One Part Of Governance}


Te Hiku's public material places its Kaitiakitanga License alongside its data-sovereignty work. That makes the licence an important artefact to examine, but not a magic label to copy into an unrelated project. A licence expresses permissions and obligations. Authority, relationships, decision forums, access processes and the capacity to enforce those terms still have to exist around it. \autocite{te-hiku-kaitiakitanga}


\index{open source}
This is where the comparison with open-source licensing becomes useful if kept narrow. An open-source licence answers questions about permissions attached to copyrighted software. A community-designed data or content licence can carry different purposes and relationships. Neither form tells an outsider that all material should be public, and neither eliminates the need to understand who had authority to issue it.


For a technical team, the design implications may include:


\begin{itemize}[leftmargin=*]
\item recording provenance and authority alongside technical metadata;
\item separating permission to access material from permission to create derivatives;
\item making later model training or commercial use a new decision rather than an assumed extension;
\item designing revocation, suspension and review paths before a breach;
\item ensuring that community decision-makers can obtain meaningful audit evidence;
\item budgeting for the people who exercise cultural and data authority.
\end{itemize}

The relevant community decides which controls fit. The engineer's role is to make those decisions enforceable and observable, not to invent cultural rules.



\subsection{Community Capability Is A Technical Outcome}


Te Hiku's public technology work includes tools for speech recognition, synthetic voice and participation in teaching computers te reo Māori. That makes capability itself part of the outcome: language technology is being developed within an organisation whose media, archives and community relationships give it a purpose beyond a benchmark score. \autocite{te-hiku-tech}


This suggests a broader way to assess a language-technology project. Model quality still matters, but so do questions such as:


\begin{itemize}[leftmargin=*]
\item Can the organisation maintain and improve the system?
\item Can it decide which future products or partners may use the work?
\item Are language experts and community authorities resourced as core contributors?
\item Does the project build durable technical capability or only extract a dataset?
\item Can users understand the boundaries attached to the output?
\end{itemize}

A project that produces an accurate model while removing future control from the community may be technically impressive and institutionally harmful. A slower project that builds authority, skills and reusable infrastructure can create more durable value.



\subsection{What Practitioners Can Carry Forward}


Four lessons travel beyond this case without pretending to reproduce it:


\begin{enumerate}[leftmargin=*]
\item \textbf{Begin with authority.} Identify who can decide the purpose, access and reuse of the material before choosing a platform.
\item \textbf{Treat governance as architecture.} Approval, provenance, access, audit and revocation are system requirements, not a policy appendix.
\item \textbf{Resource bridging work.} Language expertise, community relationships, legal understanding and technical implementation must meet inside the delivery process.
\item \textbf{Measure continuing benefit and control.} Accuracy and throughput are incomplete if the organisation cannot sustain the system or govern its future use.
\end{enumerate}

\index{CARE principles}
\index{OCAP}
Do not copy Māori terms or Te Hiku's licence into another Indigenous context as a shortcut. CARE is a global Indigenous data-governance framework; OCAP® is specifically First Nations; Te Mana Raraunga articulates Māori data sovereignty; Maiam nayri Wingara articulates principles in Australia. Each points readers back to the authority of the peoples and communities concerned. \autocite{care,ocap,te-mana-raraunga,mnw-principles}


The value of the Te Hiku example is not that it provides a universal template. It shows that language technology can be framed around community purpose, authority and stewardship from the start. The next professional step is to follow the organisation's own public material, seek the relevant expertise, and resist turning a short case study into claims the source has not made.

\section{Practice Artefact}

\index{licence}
\index{maintainer}
\index{open source}
Produce a data-sharing and stewardship memo for one dataset or open-source artefact. Classify the material, name who has authority over it, choose an access tier, identify the licence or agreement that should govern reuse, and state what benefit returns to the community or maintainers.


If the material involves Indigenous knowledge, language, community records, or cultural material, do not treat ``open'' as the default. State whose review is required before access changes and how that decision will be recorded.
